\section{Expected maximal abelian subalgebras}
\label{sec:expected}

We now combine the measured paving theorem with a separate argument for the
part of an operator orthogonal to the normalizer algebra. Throughout this
section, $M$ has separable predual, $\varphi$ is a faithful normal state, and
$A\subset M$ is a MASA contained in the centralizer $M_\varphi$. We write
$E_A$ for the $\varphi$-preserving normal conditional expectation and use
the sequence quotient of Lemma~\ref{lem:sequence-quotient}. Brackets denote
classes in that quotient; a fixed operator also denotes its constant
sequence. A partition in the quotient will always be represented by
partitions in $A$ with one fixed finite number of members.

\begin{theorem}\label{thm:state-paving}
For every $\gamma>0$ there is an integer $R_{\mathrm{st}}(\gamma)$ with the following
property. For every self-adjoint $x\in M$ with $\norm{x}\leq1$, there are
partitions $\mathcal P_j$ of $1$ in $A$, each with $R_{\mathrm{st}}(\gamma)$ members,
and self-adjoint $w_j\in M$ such that
\begin{equation}\label{exp:state-paving}
 \norm{w_j}\leq2,\qquad
 \norm{w_j}_\varphi\longrightarrow0,\qquad
 \norm{C_{\mathcal P_j}(x-w_j)-E_A(x)}\leq\gamma.
\end{equation}
The integer $R_{\mathrm{st}}(\gamma)$ is independent of $M,A,\varphi$, and $x$.
\end{theorem}

\subsection{The normalizer algebra}

Let $N$ be the von Neumann algebra generated by the groupoid normalizers
of $A$ in $M$. These are the partial isometries $v$ with $v^*v,vv^*\in A$
such that conjugation by $v$ carries $A v^*v$ onto $A vv^*$. The algebra
$N$ contains $A$. Since the modular group of $\varphi$ fixes $A$
pointwise, it carries every groupoid normalizer to another with the same
supports and the same induced corner isomorphism. Hence $N$ is globally
modular invariant. Takesaki's theorem~\cite{Takesaki72} gives a normal
$\varphi$-preserving conditional expectation $E_N:M\to N$; uniqueness of
the state-preserving expectation onto $A$ gives
\begin{equation}\label{exp:expectation-composition}
 E_A E_N=E_A.
\end{equation}

The inclusion $A\subset N$ is Cartan. For completeness, groupoid
regularity here implies regularity by unitary normalizers. Represent $A$
on a standard probability space. A normal isomorphism between two
corners of $A$ is represented by a nonsingular measurable partial
bijection. On its fixed-point set, the corresponding partial normalizer
belongs to $A$. On the complement, a countable separating family of
measurable sets partitions the domain into countably many pieces on each
of which the domain and range are disjoint. If $v$ is the restriction to
one such piece, then
\[
 v+v^*+1-v^*v-vv^*
\]
is a unitary normalizer. Multiplication by its source projection recovers
$v$, and strong summation recovers the original partial normalizer.
Also $A$ is a MASA in $N$ and $E_A|_N$ is a faithful normal expectation,
as required for the Cartan theorem.

The Feldman--Moore representation~\cite{FM77} identifies $A$ with
$L^\infty(V,\mu)$, where $\mu$ is the probability measure given by
$\varphi|_A$, and identifies $N$ with the algebra of a countable
nonsingular equivalence relation, allowing a scalar unitary cocycle.
We use the representation on square-integrable columns, with the right
base point integrated against $\mu$ and counting measure in its orbit.
An algebra element acts by left matrix multiplication on the orbit
column. In the twisted case, changing the base point changes the matrix
by diagonal unitary conjugacy, by the scalar cocycle identity. The
conditional expectation onto $A$ is its rooted diagonal. This is
$E_A|_N$: if $F:N\to A$ is any normal conditional expectation,
bimodularity gives $F(C_{\mathcal Q}(z))=F(z)$. The bounded
compression net converges ultraweakly to $E_A(z)$ by
Lemma~\ref{lem:compression-net}, so normality gives
$F(z)=F(E_A(z))=E_A(z)$.

Bounded graph partial-isometry kernels and diagonal coefficients generate
$N$. Their algebraic span, including products, is strongly dense. A
finite such expression has off-diagonal support in a finite union of
partial-bijection graphs, and therefore in an undirected graph of
uniformly bounded degree. Operator-norm bounds in the column
representation give the same matrix bounds for almost every orbit. All
models and cocycle identities below may be restricted to a common
conull invariant set for the countably many operators involved.

Fix $x=x^*$ with $\norm{x}\leq1$, and put
\begin{equation}\label{exp:decomposition}
 a_0=E_Ax,\qquad Y=E_Nx-a_0,\qquad Z=x-E_Nx.
\end{equation}
Then $x-a_0=Y+Z$, with $\norm{Y},\norm{Z}\leq2$ and $E_NZ=0$.
The measured theorem will control $Y$ through its relation matrices.
For $Z$, a unitary dilation with a diffuse left--right kernel will provide
random diagonal compressions whose moments approach a free compression law. Its norm
bound will make the spectral excess state-null, giving a second
quotient norm estimate. Common refinement combines the two bounds;
clipping the resulting block-diagonal residual then supplies the
correction in Theorem~\ref{thm:state-paving}.

\begin{lemma}\label{exp:normalizer-paving}
Let $b,\sigma>0$ and let $r$ be an integer such that
\[
 2\sqrt{2/r}+2/r<b.
\]
There are sequences of $A$-partitions with $(r+1)^2$ members whose
quotient partition $\mathcal P$ satisfies
\begin{equation}\label{exp:normalizer-bound}
 \norm{C_{\mathcal P}[Y]}\leq6(Kb+\sigma),
\end{equation}
where $K$ is the universal constant in
Theorem~\ref{thm:measured-paving}.
\end{lemma}

\begin{proof}
Bounded Kaplansky approximation, followed by norm approximation inside
the generating algebra, gives self-adjoint finite graph expressions
converging strongly-star to $Y$, with norms at most $3$. Subtracting
their $A$-expectations makes their diagonals zero and increases the norm
bound to $6$; state-norm contractivity of $E_A$ preserves convergence.
For each such expression, cut down by the base set on which every
incident edge in its finite-degree enclosing graph has cocycle ratio
between $J^{-1}$ and $J$. These sets increase to full measure as
$J\to\infty$, and the cuts converge strongly-star to the expression.
A diagonal choice therefore gives self-adjoint $Y_j$ with
\[
 E_AY_j=0,\qquad \norm{Y_j}\leq6,\qquad
 \norm{Y_j-Y}_\varphi^\#\longrightarrow0,
\]
whose orbit matrices have bounded-degree support and bounded cocycle
ratios, with the bounds allowed to depend on $j$.

Apply Theorem~\ref{thm:measured-paving} to $Y_j/6$ and separately to
$-Y_j/6$, using state-error tolerances tending to zero. The rooted
diagonal formula for the state says that the positive part above
$Kb+\sigma$ has $\varphi$-norm tending to zero in each case. These
positive parts are self-adjoint and uniformly bounded, so they belong
to the null ideal of Lemma~\ref{lem:sequence-quotient}. The sequence
$Y_j-Y$ belongs to that ideal as well. The resulting two quotient
partitions, each of size $r+1$, therefore give respectively
\[
 C_{\mathcal P^+}[Y/6]\leq(Kb+\sigma)1,
 \qquad
 C_{\mathcal P^-}[Y/6]\geq-(Kb+\sigma)1.
\]
Take their common refinement. All partition projections commute, and
compression is positive and unital, so both scalar bounds hold after
refinement. This proves the lemma. Further finite $A$-refinement
preserves the bound.
\end{proof}

We have completed the normalizer part. The remaining component has the
following target, whose proof is given after the kernel and free-moment
lemmas below.

\begin{proposition}\label{prop:off-normalizer}
For every $n\geq2$, there are sequences of $n$-partitions in $A$
whose quotient partition $\mathcal P$ satisfies
\begin{equation}\label{exp:off-normalizer-bound}
 \norm{C_{\mathcal P}[Z]}\leq\frac4{\sqrt n}.
\end{equation}
\end{proposition}

\subsection{The diffuse kernel outside the normalizer algebra}
We retain the faithful normal state $\varphi$, with $A\subset M_\varphi$,
and the algebra $N$ generated by the groupoid normalizers of $A$.
Use a standard probability model $A=L^\infty(V,\mu)$ for
$\varphi|_A$.  We use the standard-form realization $L^2(M)$, equivalently
Haagerup $L^2(M)$, and write $\varphi^{1/2}$ for the positive vector
implementing $\varphi$.  In particular,
$a\varphi^{1/2}=\varphi^{1/2}a$ for $a\in A$.

For $X\in M$, the commuting left and right actions of $A$ give a finite
joint spectral measure $\nu_X$ for $X\varphi^{1/2}$ on $V\times V$.
Our first coordinate is the right coordinate: for Borel sets $B,C\subset V$,
\begin{equation}\label{exp:joint-measure}
 \nu_X(B\times C)
 =\big\|1_CX\varphi^{1/2}1_B\big\|_2^2
 =\int_B E_A(X^*1_CX)(v)\,d\mu(v).
\end{equation}
Both marginal measures are absolutely continuous with respect to $\mu$.
The right marginal has density $E_A(X^*X)\leq\|X\|^2$; the left marginal
is absolutely continuous because $a\mapsto\varphi(X^*aX)$ is normal on
$A$.  Disintegration on standard Borel spaces therefore gives a measurable
kernel $k_X(v,dw)$ such that
\begin{equation}\label{exp:kernel}
 \nu_X(dv,dw)=d\mu(v)k_X(v,dw),\qquad
 E_A(X^*aX)(v)=\int_V a(w)\,k_X(v,dw),\qquad
 k_X(v,V)\leq\|X\|^2
\end{equation}
almost everywhere, for $a\in A$.  As usual, individual function identities
are understood modulo null sets.  The joint spectral resolution can be
constructed from a bounded real coordinate generating the standard model
and the commuting left and right spectral resolutions of that coordinate.

\begin{lemma}\label{lem:atomless}
 If $E_N(X)=0$, then $k_X(v,\cdot)$ has no atoms for almost every $v\in V$.
\end{lemma}

\begin{proof}
Choose Borel representatives for the kernel.  The set
\[
 \mathcal T=\{(v,w):k_X(v,\{w\})>0\}
\]
is Borel and has countable sections in the second coordinate.  For example,
joint measurability of the singleton mass follows by using shrinking
neighborhoods in a compatible separable metric on the standard Borel
space.  The countable-section uniformization theorem expresses
$\mathcal T$ as a countable union of graphs of partial Borel maps.
If the conclusion fails, one such map $F:S\to V$ has a domain of positive
$\mu$-measure, with
\[
 g(v):=k_X(v,\{F(v)\})>0\quad(v\in S).
\]
There is no assumption that $F$ is injective.  For a $\mu$-null Borel set
$B\subset V$, absolute continuity of the left marginal gives
\[
 \int_{S\cap F^{-1}(B)}g(v)\,d\mu(v)
 \leq\nu_X(V\times B)=0.
\]
Thus $F$ pulls back null sets to null sets.  With $e=1_S\in A$, composition
with $F$ consequently defines a normal $*$-homomorphism
\[
 \theta:A\longrightarrow Ae,\qquad
 \theta(a)(v)=a(F(v))\quad(v\in S),\qquad \theta(1)=e.
\]
The graph will intertwine $A$ with the possibly proper subalgebra
$\theta(A)$ of $Ae$; its polar partial isometry need not yet normalize
$A$. We will use an expected type I intermediate algebra to split the
right support into corners where these two abelian algebras coincide.

Apply the joint spectral projection of the graph of $F$ to
$X\varphi^{1/2}$ and call the resulting vector $\zeta$.  It is nonzero and
satisfies
\begin{equation}\label{exp:graph-intertwining}
 \zeta e=\zeta,\qquad a\zeta=\zeta\theta(a)\quad(a\in A).
\end{equation}
Its right spectral density on $S$ is $g$.  In particular, $e$ is its
smallest right supporting projection among the projections of $A$:
if $f\in A$ is a projection and $\zeta f=\zeta$, then $e\leq f$.

Take the $L^2$ polar decomposition
$\zeta=u h$, where $h=|\zeta|\in L^2(M)_+$ and $u\in M$ is a partial
isometry.  Put $p=u^*u\leq e$ and $q=uu^*$.  We use here the usual
polar decomposition and support properties of standard-form
noncommutative $L^2$; see, for example, \cite[Lemma~11.12 and Theorem~11.30]{Hiai20}.
For a unitary $a\in A$, the operator $b=\theta(a)+1-e$ is unitary and
\eqref{exp:graph-intertwining} reads $a\zeta=\zeta b$.
Taking absolute values gives $h=b^*hb$.  Uniqueness of the polar
decomposition then gives $au=ub=u\theta(a)$.  Consequently,
\begin{equation}\label{exp:polar-intertwining}
 h\theta(a)=\theta(a)h,\qquad
 p\theta(a)=\theta(a)p,\qquad
 au=u\theta(a)\quad(a\in A),\qquad q\in A.
\end{equation}
The identities for all $a$ follow by linearity from unitaries.  For the
last assertion, the identities imply that $q$ commutes with $A$, and $A$
is maximal abelian.

Consider the von Neumann algebra
\[
 D=\theta(A)'\cap eMe.
\]
It contains $p$ and $Ae$, and $Ae$ is a MASA in $D$.
Moreover $\theta(A)\subset Z(D)$.  If $d\in pDp$, then $udu^*$ commutes
with $Aq$, by \eqref{exp:polar-intertwining}, and hence belongs to $Aq$.
On the other hand,
\[
 u\theta(a)p u^*=aq\quad(a\in A).
\]
Conjugation by $u$ therefore identifies $pDp$ with $Aq$, and
\begin{equation}\label{exp:abelian-corner}
 pDp=\theta(A)p.
\end{equation}
Let $c$ be the central support of $p$ in $D$.  Since
$Z(D)\subset Ae$, the identity $pc=p$ gives $\zeta c=\zeta$; the minimal
right $A$-support property gives $c=e$.  Thus $D$ has an abelian projection
with full central support and is type I.  Finally, if $d\in Z(D)$, then
\eqref{exp:abelian-corner} gives $dp=\theta(a)p$ for some $a\in A$.
The central element $d-\theta(a)$ annihilates a projection of full central
support and is zero.  We have proved
\begin{equation}\label{exp:center}
 Z(D)=\theta(A).
\end{equation}

We next need a partition of $e$ by projections $e_j\in Ae$ such that
\begin{equation}\label{exp:branch-partition}
 (Ae)e_j=\theta(A)e_j.
\end{equation}
The restriction of $E_A$ to $D$ is a normal conditional expectation onto
$Ae$.  The type I MASA theorem of Akemann and Sherman
\cite[Theorem~4.1]{AS12} supplies a partition of $e$ by projections in
$Ae$ that are abelian in $D$.  There are only countably many nonzero
members, because $D$ has the faithful normal state
$\varphi|_D/\varphi(e)$.  For an abelian projection in a type I algebra,
\[
 e_jDe_j=Z(D)e_j.
\]
This follows, for instance, by writing the type I algebra as a direct
integral: an abelian projection has fiber rank at most one.  Together
with \eqref{exp:center} and $\theta(A)\subset Ae\subset D$, this proves
\eqref{exp:branch-partition}.

Set $\eta_j=\zeta e_j$ and take its polar decomposition
$\eta_j=u_j h_j$.  Its right support is at most $e_j$.  Repeating
\eqref{exp:polar-intertwining} shows that $h_j$ commutes with $\theta(A)$
and that
\[
 au_j=u_j\theta(a)\quad(a\in A).
\]
By \eqref{exp:branch-partition} and the support condition, $h_j$ commutes
with all of $A$.  Therefore the right support $p_j=u_j^*u_j$ belongs to
$A$; the left support $q_j=u_ju_j^*$ belongs to $A$ by the same polar
argument as before.  In addition,
\[
 \theta(A)p_j=Ap_j,\qquad
 u_j(Ap_j)u_j^*=Aq_j.
\]
Thus each nonzero $u_j$ is a groupoid normalizer and belongs to $N$.

Let $\psi_j$ be the normal positive functional whose positive
implementing vector is $h_j$.  Since $h_j$ commutes with $A$, this
functional centralizes $A$.  In particular
$\psi_j\circ C_{\mathcal Q}=\psi_j$ for every finite partition
$\mathcal Q$ in $A$.  Lemma~\ref{lem:compression-net}, normality, and
boundedness of the compression net imply
\[
 \psi_j=\psi_j\circ E_A.
\]
Write $f_j=d(\psi_j|_A)/d\mu\in L^1(V,\mu)_+$ and
$f_{j,m}=\min(f_j,m)$.  The positive vectors
$f_{j,m}^{1/2}\varphi^{1/2}$ belong to $A\varphi^{1/2}$ and converge in
$L^2$; indeed their squared pairwise distances are
\[
 \int_V\big|f_{j,m}^{1/2}-f_{j,n}^{1/2}\big|^2\,d\mu.
\]
Their normal positive functionals are
$x\mapsto\int f_{j,m}E_A(x)\,d\mu$ and converge in norm to $\psi_j$.
The limit is therefore the unique positive implementing vector $h_j$.
It follows that
\[
 h_j\in\overline{A\varphi^{1/2}},\qquad
 \eta_j\in\mathcal H_N:=\overline{N\varphi^{1/2}}.
\]
The right projections $e_j$ are orthogonal and sum to $e$, so
$\zeta=\sum_j\zeta e_j\in\mathcal H_N$, with convergence in $L^2$.

The closed space $\mathcal H_N$ reduces both the left and right actions
of $A$: for the right action, use
$n\varphi^{1/2}a=na\varphi^{1/2}$.  It therefore reduces their joint
spectral projections, including the graph projection defining $\zeta$.
But $E_N(X)=0$ gives, for $n\in N$,
\[
 \varphi(n^*X)=\varphi(n^*E_N(X))=0,
\]
so $X\varphi^{1/2}\perp\mathcal H_N$.  Thus
$\zeta\perp\mathcal H_N$, contradicting $\zeta\ne0$ and the
membership just proved.
\end{proof}

\paragraph{Removing atomic summands.}
For the remaining treatment of $Z=x-E_N(x)$ we may assume that $A$ is
diffuse.  Here are the details of this reduction.  If $f$ is an atom of
$A$, then $Af=\mathbb C f$ is a MASA in $fMf$, so $fMf=\mathbb C f$.
Thus $f$ is minimal in $M$.  Let $z_f$ be its central support in $M$.
The algebra $z_fM$ is type I, since $f$ is an abelian projection with
full central support there.  Compression by $f$ is faithful on its
center and has scalar range, so $z_fM$ is a type I factor.

The algebra $Az_f$ is a MASA in this factor and has the normal
expectation $E_A|_{z_fM}$.  Again by \cite[Theorem~4.1]{AS12}, it has a
partition of $z_f$ into rank-one projections.  In particular it is
purely atomic.  Partial isometries joining any two of these rank-one
projections normalize their $A$-corners.  A compatible system of these
matrix units generates $z_fM$, and hence
\[
 z_fM\subset N.
\]
Distinct central supports of atoms are either equal or orthogonal,
because each supports a factor summand.  Their sum $z_{\mathrm{at}}$ is
a central projection, belongs to $A$, and satisfies
$z_{\mathrm{at}}M\subset N$.  Consequently
\[
 z_{\mathrm{at}}Z=0.
\]
On the complementary central summand, $A(1-z_{\mathrm{at}})$ has no
atoms.  The algebra generated there by its groupoid normalizers is
$N(1-z_{\mathrm{at}})$: corner normalizers extend by zero, and global
normalizers restrict to this central summand.  We may use the normalized
restriction of $\varphi$ on the nonzero complementary summand.  This
changes state seminorms only by a fixed scalar factor.
A partition obtained there extends to $M$ with the same number of cells
by adding $z_{\mathrm{at}}$ to one cell; all estimates for $Z$ are
preserved because $Z$ vanishes on $z_{\mathrm{at}}$.


\subsection{Random diagonal unitaries and free moments}

It remains to compress the diffuse summand of $Z$.
The next two lemmas supply the free law and its norm bound for the
matrix dilation used to prove Proposition~\ref{prop:off-normalizer}.
The state in the probabilistic argument need not be tracial.

\begin{lemma}\label{lem:random-freeness}
Let $\mathcal M$ have separable predual, let $\phi_0$ be a faithful
normal state, and let $D\subset\mathcal M_{\phi_0}$ be a diffuse MASA.
Suppose that $U\in\mathcal M$ is a self-adjoint unitary with
$\phi_0(U)=0$, and that the normal positive map
\[
 \mathcal K(a)=E_D(UaU),\qquad a\in D,
\]
has an almost everywhere diffuse kernel on a standard probability
model of $(D,\phi_0|_D)$. For every integer $n\geq2$ there are unitaries
$w_l\in D$, with $w_l^n=1$, such that the joint star moments of
$(w_l,U)$ converge to those of free unitaries $(w,U_{\mathrm f})$,
where $w$ is uniform on the $n$th roots of unity and $U_{\mathrm f}$
is uniform on $\{-1,1\}$.
\end{lemma}

\begin{proof}
Choose an increasing sequence of finite partitions $\mathcal Q$ in
$D$ which generates $D$, separates points of a conull standard model,
and satisfies
\begin{equation}\label{exp:small-cells}
 \max_{q\in\mathcal Q}\phi_0(q)\longrightarrow0.
\end{equation}
Such a sequence is obtained by adjoining a countable separating family
and, at each stage, a sufficiently small-mass finite refinement; the
latter exists by diffuseness. Give its cells independent labels $Z_q$
uniform on the $n$th roots of unity, and put $w_{\mathcal Q}=\sum_q Z_q q$.

We first make the averaged reduced moments vanish by estimating return
probabilities for diffuse kernels. We then control their fluctuations
using the small masses of the partition cells, so that one deterministic
labeling works for each prescribed finite list of moments.

\paragraph{Averaged words and a repeated-cell compression.}
We first prove that the label expectation tends to zero for each scalar
\begin{equation}\label{exp:reduced-test-word}
 F_{\mathcal Q}
 =\phi_0(x_0 w_{\mathcal Q}^{k_1}x_1\cdots
                 w_{\mathcal Q}^{k_h}x_h),
 \qquad h\geq1,\quad 1\leq k_j<n,
\end{equation}
where $x_j=U$ for $1\leq j<h$ and each endpoint $x_0,x_h$ is either
$1$ or $U$. Expand over cell assignments $q_1,\ldots,q_h$.
An assignment survives label averaging only if, in every block of equal
cells, the sum of its exponents is zero modulo $n$. Consequently every
block of its exact equality pattern has size at least two. For any
surviving exact pattern, inclusion-exclusion expresses its word sum as
a finite linear combination of sums
\begin{equation}\label{exp:partition-word-sum}
 S_\pi
 =\sum_{\substack{q_1,\ldots,q_h\in\mathcal Q\\
                    q_i=q_j\text{ if }i,j\text{ are in one block of }\pi}}
      \phi_0(x_0q_1x_1\cdots q_hx_h),
\end{equation}
where every block of $\pi$ has at least two members. There is no
distinctness requirement between different blocks in
\eqref{exp:partition-word-sum}. It is enough to show $S_\pi\to0$ for
each such $\pi$.

Choose positions $c<d$ in one block whose distance $d-c$ is minimal
among all pairs in a common block. The positions strictly between them
belong to distinct other blocks. We impose the equalities of $\pi$
using auxiliary Haar phases. For a block other than the selected block,
place independently sampled cellwise Haar diagonal unitaries at all
but one of its occurrences, and their inverse product at the remaining
occurrence. Haar averaging forces precisely the equality of the cells
at those occurrences. In the selected block, use the same projection
$q$ explicitly at $c$ and $d$ and sum over $q$; put independent Haar
diagonals at every other occurrence of that block, with their inverse
product additionally at $c$. All initial draws are independent.
Expanding and averaging now gives exactly $S_\pi$ as the auxiliary
expectation of
\begin{equation}\label{exp:haar-compression}
 \phi_0\bigl(\mathrm{PRE}\,C_{\mathcal Q}(Y_{\mathcal Q})\,
                    \mathrm{POST}\bigr),\qquad
 Y_{\mathcal Q}=U d_{c+1}U\cdots d_{d-1}U.
\end{equation}
When $d=c+1$, the last expression is $U$. The extra inverse product at
$c$ can be moved before its projection, since both lie in $D$, and is
included in $\mathrm{PRE}$. The factors outside the compression are
products of fixed contractions $x_j$ and diagonal unitaries. Each
interior $d_j$ is a cellwise Haar diagonal: this is also true if it is
the inverse product placed at the last occurrence of its block. The
interior diagonals are mutually independent because their positions
belong to different blocks. Independence from $\mathrm{PRE}$ and
$\mathrm{POST}$ is not asserted or needed.

Their independence and orthogonality of the outside projections give
\begin{equation}\label{exp:haar-square-sum}
 \E\norm{C_{\mathcal Q}(Y_{\mathcal Q})}_{\phi_0}^2
 =\sum_{q,q_{c+1},\ldots,q_{d-1}\in\mathcal Q}
   \norm{qUq_{c+1}U\cdots q_{d-1}Uq}_{\phi_0}^2.
\end{equation}
There are no middle cells when $d=c+1$. Fix a coarser partition
$\mathcal E$ from the generating sequence. When $\mathcal Q$ refines
$\mathcal E$, contractivity of $C_{\mathcal Q}$ in the state norm,
applied with each middle cell tuple fixed, bounds the right side by
\begin{equation}\label{exp:coarse-square-sum}
 \sum_{p\in\mathcal E}
 \sum_{q_{c+1},\ldots,q_{d-1}\in\mathcal Q}
   \norm{pUq_{c+1}U\cdots q_{d-1}Up}_{\phi_0}^2.
\end{equation}

\paragraph{Diffuse kernel return probabilities.}
Here is the limit of this expression without a tracial assumption.
Put $m=d-c$. For each $p\in\mathcal E$ define
\[
 B_{1,\mathcal Q}(p)=UpU,\qquad
 B_{j+1,\mathcal Q}(p)=U C_{\mathcal Q}(B_{j,\mathcal Q}(p))U.
\]
Expanding the squares in \eqref{exp:coarse-square-sum} and summing the
middle cells identifies that expression with
$\sum_{p\in\mathcal E}\phi_0(pB_{m,\mathcal Q}(p)p)$.
All $B_{j,\mathcal Q}(p)$ are positive contractions. By
Lemma~\ref{lem:compression-net}, $C_{\mathcal Q}(t)\to E_D(t)$
strongly-star on each fixed bounded element. Its contractivity for
both state seminorms also controls a bounded, varying, strong-star
null error. Bounded strong-star continuity of multiplication by the
fixed $U$ therefore proves recursively that $B_{j,\mathcal Q}(p)$
converges strongly-star to $B_j(p)$, where
\[
 B_1(p)=UpU,\qquad B_{j+1}(p)=U E_D(B_j(p))U.
\]
Thus $E_D(B_j(p))=\mathcal K^j(p)$. If $\mu_0$ denotes the probability
measure on $D$, state preservation and bimodularity of $E_D$ show that
the limit of \eqref{exp:coarse-square-sum} is
\begin{equation}\label{exp:kernel-return}
 \sum_{p\in\mathcal E}\phi_0(p\mathcal K^m(p)p)
 =\sum_{p\in\mathcal E}\int_p(\mathcal K^m1_p)(v)\,d\mu_0(v).
\end{equation}

The positive normal map $\mathcal K$ is unital. Its kernel is
nonsingular in the integrated sense: a fixed $\mu_0$-null set receives
zero mass for almost every starting point. This follows from
normality, since its indicator is zero in $D$. In particular the
kernel avoids, for almost every starting point, the fixed exceptional
null set on which its successor kernel is not diffuse. Integrating
diffuse successor kernels consequently produces a diffuse kernel.
Induction shows that each fixed iterate $\mathcal K^m$ has an almost
everywhere diffuse kernel. As $\mathcal E$ increases along the
separating sequence, the cell containing $v$ decreases to $\{v\}$ on
the chosen conull model. Its return mass in
\eqref{exp:kernel-return} therefore tends to zero for almost every
$v$. It is bounded by $\mathcal K^m1(v)=1$. Dominated convergence
proves that \eqref{exp:kernel-return} tends to zero. First fixing
$\mathcal E$ and letting $\mathcal Q$ refine, and then letting
$\mathcal E$ refine, proves
\begin{equation}\label{exp:haar-square-zero}
 \E\norm{C_{\mathcal Q}(Y_{\mathcal Q})}_{\phi_0}^2\longrightarrow0.
\end{equation}

\paragraph{Nontracial state estimates.}
We next explain the state estimates that turn this into scalar moment
convergence. For an entire analytic element $a$ for the modular group,
the standard-form right-action identity gives
\begin{equation}\label{exp:analytic-right}
 \norm{ta}_{\phi_0}
 \leq\norm{\sigma_{i/2}^{\phi_0}(a)}\norm{t}_{\phi_0}.
\end{equation}
See \cite[Lemma~8.12, Equation~(8.4)]{Hiai20}, applied to $a^*$.
Entire analytic elements approximate each fixed element boundedly in
the state norm, for example by Gaussian modular smoothing
\cite[Lemma~2.13]{Hiai20}. A centralizer unitary has right-multiplication bound
$1$. In any scalar word consisting of fixed bounded factors and
varying centralizer unitaries, replace the fixed factors from right to
left by bounded analytic approximants. At a replacement step all
factors to the right already have uniformly bounded right actions,
and all factors to the left have uniformly bounded operator norms.
The error in the scalar state value is therefore at most a fixed
constant times the state-norm error of the factor being replaced.
Choosing the approximants in that order makes the total error
arbitrarily small, uniformly over the varying diagonals and
partitions.

Apply this procedure only to the fixed factors in $\mathrm{POST}$
in \eqref{exp:haar-compression}. The compression and all remaining
left factors are contractions. With the approximants fixed, the
absolute scalar value is bounded by
$C\norm{C_{\mathcal Q}(Y_{\mathcal Q})}_{\phi_0}$, where $C$ is
independent of every phase and of $\mathcal Q$. Cauchy--Schwarz in
the auxiliary probability space and \eqref{exp:haar-square-zero}
make its phase expectation tend to zero. The uniformly arbitrarily
small replacement error proves $S_\pi\to0$. Hence
\begin{equation}\label{exp:word-expectation-zero}
 \E F_{\mathcal Q}\longrightarrow0.
\end{equation}

\paragraph{Concentration and deterministic labels.}
It remains to choose deterministic labels simultaneously for the
moments. Replace all fixed factors in
\eqref{exp:reduced-test-word} analytically, again from right to left,
and call the resulting scalar $G_{\mathcal Q}$. It approximates
$F_{\mathcal Q}$ uniformly to any prescribed accuracy. Changing one
label $Z_q$ changes a power of $w_{\mathcal Q}$ by a scalar of
absolute value at most $2$ times $q$. Telescoping over the finitely
many occurrences, each difference is thus a bounded scalar times
$\phi_0(\mathrm{PRE}\,q\,\mathrm{POST})$. Cauchy--Schwarz gives
\begin{align}\label{exp:two-sided-influence}
 \abs{\phi_0(\mathrm{PRE}\,q\,\mathrm{POST})}
 &\leq\norm{q\mathrm{PRE}^*}_{\phi_0}
        \norm{q\mathrm{POST}}_{\phi_0}\notag\\
 &\leq C_{\mathrm{word}}\phi_0(q).
\end{align}
The second line follows from \eqref{exp:analytic-right} for both
products, including their analytic adjoints. Their modular norm
bounds are uniform over all other labels, since the diagonal
unitaries are in the centralizer. Independent bounded differences
(or the martingale variance sum obtained by revealing one label at
a time) now give
\[
 \E\abs{G_{\mathcal Q}-\E G_{\mathcal Q}}^2
 \leq C'_{\mathrm{word}}\sum_{q\in\mathcal Q}\phi_0(q)^2
 \longrightarrow0
\]
by \eqref{exp:small-cells}. The uniform approximation shows that
$F_{\mathcal Q}-\E F_{\mathcal Q}\to0$ in probability. Together with
\eqref{exp:word-expectation-zero}, this gives
$F_{\mathcal Q}\to0$ in probability.

There are only countably many test words. For successively larger
finite lists, convergence in probability gives a partition far
enough along the sequence and a deterministic labeling for which
every word in the list is as small as prescribed. A diagonal choice
gives $w_l$ with all those moments tending to zero. The tests include
the nontrivial single powers of $w_l$, so its limiting distribution
is uniform on the $n$th roots. Also $U^2=1$ and $\phi_0(U)=0$ give
the symmetric two-point distribution. Centered elements of the two
limiting generated algebras are linear combinations of the tested
powers and of $U$, respectively; all alternating centered moments
are precisely the tests above. This establishes the asserted joint
free law.
\end{proof}

\begin{lemma}\label{lem:free-compression}
Let $w$ and $U_{\mathrm f}$ have the free law in
Lemma~\ref{lem:random-freeness}, in its reduced free-product
representation. If $p_1,\ldots,p_n$ are the spectral projections of
$w$, then
\begin{equation}\label{exp:free-norm-bound}
 \norm{\sum_{j=1}^n p_j U_{\mathrm f}p_j}\leq\frac2{\sqrt n}.
\end{equation}
\end{lemma}

\begin{proof}
Let $E_0$ be the Hilbert-space projection onto the vacuum and all
reduced words not starting in the centered part of the $w$-algebra.
For a spectral projection $p$ of $w$, put $t=1/n$. On this subspace,
left multiplication by $p$ is the sum of the scalar action $t$ and
creation by $p-t1$. The two resulting subspaces are orthogonal, and
$\norm{p-t1}_2^2=t(1-t)$; hence
\[
 \norm{pE_0}=\norm{E_0p}=\sqrt{t}.
\]
Since $U_{\mathrm f}$ is centered, it sends a word starting in the
$w$-centered space to a word starting in the other factor. Therefore
$U_{\mathrm f}(1-E_0)=E_0U_{\mathrm f}(1-E_0)$, and
\[
 pU_{\mathrm f}p
 =pU_{\mathrm f}E_0p+pE_0U_{\mathrm f}(1-E_0)p.
\]
Unitarity gives norm at most $2\sqrt t$. The sum of the compressions
is block diagonal for the mutually orthogonal $p_j$, so its norm is
the largest block norm.
\end{proof}

\subsection{A free-product matrix dilation}

\begin{proof}[Proof of Proposition~\ref{prop:off-normalizer}]
By the atomic reduction following Lemma~\ref{lem:atomless}, it is
enough to work on the central summand where $A$ is diffuse. Normalize
the state on that summand and retain the same notation. Put $z=Z/2$,
so $z$ is a self-adjoint contraction and $E_Nz=0$.

The ordinary two-by-two unitary dilation need not have a diffuse
kernel over $A\oplus A$: already at $z=0$, its identity off-diagonal
entry gives an atomic cross-kernel. We replace the second diagonal
algebra by a free diffuse abelian algebra, turning the cross terms
into normal functionals on diffuse spaces. This is the symmetry form
of the free-product projection dilation in
\cite[proof of Theorem~5.1]{PV15}: with $y=(1+z)/2$, the symmetry
below is twice that projection minus the identity. The kernel and
moment lemmas above supply the required freeness in the present
nontracial setting.

Let $(B_0,\tau_B)$ be a diffuse standard separable abelian probability
algebra and form the reduced von Neumann free product
\[
 (P,\psi)=(M,\varphi)*(B_0,\tau_B).
\]
The free-product state is faithful and normal, the two factor
expectations preserve it, and its Hilbert-space representation is
the reduced-word representation; see, for example,~\cite{HU}.
Both $A$ and $B_0$ remain MASAs in $P$. We give the argument to
make its nontracial applicability explicit. There are unitaries
$v_k\in A$ tending weakly to zero, since $A$ is diffuse. They
centralize $\psi$, by the expectation onto $M$. On
$L^2(P,\psi)\ominus L^2(M,\varphi)$, a dense family consists of
vectors of words
\[
 x_0 b_1 x_1\cdots b_hx_h,\qquad h\geq1,
\]
where $b_i\in B_0$ are centered, the internal $x_i\in M$ are
centered, and the endpoints are arbitrary. To compute a matrix
coefficient of conjugation by $v_k$, take another such word with
initial endpoint $y_0$, and split the middle product
$y_0^*v_kx_0$ into its centered part and its scalar mean. The
centered part gives a reduced word with zero expectation onto $M$.
The scalar part is $\varphi(y_0^*v_kx_0)$ times an expression with
a fixed norm bound and tends to zero. The additional right factor
$v_k^*$ in the conjugation does not change that norm bound.
Conjugation is unitary on the state Hilbert space; density proves
that it tends weakly to zero on the indicated complement. If an
element of $P$ commutes with $A$, its vector after subtracting its
$M$-expectation lies in this complement and is fixed by all these
conjugations. It must be zero. The element belongs to $M$, and
then to $A$. Interchanging the factors proves the assertion for
$B_0$.

In $\widetilde M=\operatorname{Mat}_2(P)$ use the state
$\widetilde\varphi=\tr_2\otimes\psi$, with normalized matrix trace.
The diagonal algebra $D=A\oplus B_0$ is a diffuse MASA in its
centralizer. Define
\begin{equation}\label{exp:unitary-dilation}
 U=\begin{pmatrix}z&s\\s&-z\end{pmatrix},\qquad
 s=(1-z^2)^{1/2}.
\end{equation}
Since $z$ commutes with $s$, this is a self-adjoint unitary.
Furthermore $E_D(U)=0$. For $a\in A$ and $b\in B_0$,
direct free-product expectation gives
\begin{equation}\label{exp:dilation-kernel}
 E_D\bigl(U(a\oplus b)U\bigr)
 =\bigl(E_A(zaz)+\tau_B(b)E_A(s^2)\bigr)
  \oplus
  \bigl(\varphi(sas)1+\tau_B(b)\varphi(z^2)1\bigr).
\end{equation}
In particular, $E_{B_0}(zbz)=\tau_B(b)\varphi(z^2)1$ because
$\varphi(z)=0$. Lemma~\ref{lem:atomless} makes the first kernel
$a\mapsto E_A(zaz)$ diffuse. The other kernels in
\eqref{exp:dilation-kernel} are diffuse because the base algebras
are diffuse and the displayed functionals are normal. Thus the
kernel of $a\mapsto E_D(UaU)$ is almost everywhere diffuse.

Apply Lemma~\ref{lem:random-freeness} with $\phi_0=
\widetilde\varphi$. Let $P_{j,l}\in D$, $1\leq j\leq n$, be the
spectral projections of the resulting finite-order $w_l$, padded
with zero projections when necessary, and set
\[
 S_l=\sum_{j=1}^n P_{j,l}UP_{j,l}.
\]
Each projection is a polynomial in $w_l$ and $w_l^*$.
Consequently the scalar moments of the self-adjoint contractions
$S_l$ tend to the moments of
$S_{\mathrm f}=\sum_jp_jU_{\mathrm f}p_j$ in the reduced free law.
Let $c=2/\sqrt n$. Lemma~\ref{lem:free-compression} gives
$\norm{S_{\mathrm f}}\leq c$. Uniform polynomial approximation
on $[-1,1]$, applied to the continuous function
$t\mapsto(\abs{t}-c)_+^2$, gives
\begin{equation}\label{exp:excess-null}
 \norm{(\abs{S_l}-c)_+}_{\widetilde\varphi}\longrightarrow0.
\end{equation}
This passage uses moments to obtain a state-norm estimate, not
operator-norm convergence in $\widetilde M$.

The sequence $(S_l)$ belongs to the multiplier algebra of
Lemma~\ref{lem:sequence-quotient}: $U$ is fixed and each
$(P_{j,l})_l$ is a bounded centralizer sequence. Its excess in
\eqref{exp:excess-null} is self-adjoint, so it belongs to the
strong-star null ideal. Continuous functional calculus in the
quotient gives $\norm{[S_l]}\leq c$.

Compress by the fixed first matrix projection $e_{11}$. Write
$P_{j,l}=p_{j,l}\oplus b_{j,l}$, so the $(p_{j,l})_{j=1}^n$ form
a partition in $A$. The compressed sequence is
\[
 e_{11}S_le_{11}=\sum_{j=1}^n p_{j,l}z p_{j,l}
\]
in the first corner. Its quotient norm upstairs is at most $c$.
Functional calculus there makes its absolute-value excess above
$c$ strong-star null. These compressed elements belong to $M$,
and the state on their corner restricts to $\varphi/2$; hence the
same excess is strong-star null for $(M,\varphi)$. The compressed
sequence is also a multiplier sequence for $M$, since its
partition projections centralize $\varphi$. Its original quotient
norm is therefore at most $c$. Multiplying by $2$ proves
\eqref{exp:off-normalizer-bound} on the diffuse summand. Extending
the same colors over the central summand where $Z=0$ proves the
full assertion.
\end{proof}

\subsection{Completion of the state-preserving paving argument}

\begin{proof}[Proof of Theorem~\ref{thm:state-paving}]
Choose
\[
 b=\sigma=\frac{\gamma}{24(K+1)},
\]
choose $r$ with $2\sqrt{2/r}+2/r<b$, and choose $n\geq2$ with
$4/\sqrt n\leq\gamma/2$. Use Lemma~\ref{exp:normalizer-paving}
and Proposition~\ref{prop:off-normalizer}, and take the common
refinement of their partition sequences. It has
$R_{\mathrm{st}}(\gamma)=(r+1)^2n$ members after padding. Finite refinement
preserves each quotient bound, so for
\[
 T_j^0=C_{\mathcal P_j}(x-E_Ax)
\]
we have $\norm{[T_j^0]}\leq\gamma$, since the two contributions
are at most $\gamma/4$ and $\gamma/2$.

Let $\operatorname{clip}_\gamma(t)=\max(-\gamma,
\min(t,\gamma))$ for real $t$, and define
\[
 w_j=T_j^0-\operatorname{clip}_\gamma(T_j^0).
\]
The quotient norm bound and continuous functional calculus show
that $(w_j)$ has zero quotient image. It is self-adjoint, hence
$\norm{w_j}_\varphi\to0$. Also $\norm{T_j^0}\leq2$, so
$\norm{w_j}\leq2$. The operator $T_j^0$, and therefore $w_j$,
is block diagonal for $\mathcal P_j$. Since $E_Ax$ belongs to
$A$, this gives
\[
 C_{\mathcal P_j}(x-w_j)-E_Ax
 =T_j^0-w_j=\operatorname{clip}_\gamma(T_j^0),
\]
of norm at most $\gamma$. These are all the assertions of
\eqref{exp:state-paving}.
\end{proof}
